% Français 9115, batch 38 — Gallica views 741–750, the last ten views of the volume.
% Pass: Opus 5.5 (claude-opus-5-5), 2026-09-22 — first pass, unchecked against the views by a human.
%
% What the ten views hold. Not Sophie Germain's own mathematics, and not
% « Manuscript C »: the volume ends on a run of reading notes — French
% extracts, partly translated and partly summarised, from the Mémoires
% (Miscellanea) of the Turin Academy, volumes 3 and 4, with remarks of the
% copyist's own in square brackets. The hand is a regular, slanted fair
% hand, lightly corrected between the lines; whether it is hers is not
% settled by anything on the leaves. Nine views are written; v. 750 is a
% blank leaf. The leaves are single sheets on guards, written on both
% sides: the ink number stands on the recto (even view), the verso (odd
% view) carries none, and what shows at its top left is the recto's
% figure read through the paper. No microfilm target, no repeated
% exposure. One run of extracts, begun on the leaf of v. 740 (outside
% this batch), fills v. 741–749; the extracts change author and memoir
% at the headings noted below.
%
%   v. 741  (verso of ink 368, v. 740) Euler, « Eclaircissemens sur le
%           mouvement des cordes vibrantes », the title given on v. 740,
%           « 3 V. des Memoires de Turin »: discontinuous initial figures;
%           arbitrary functions entering the analysis of several
%           variables; y = Γ:(x+ct) + Δ:(x−ct); strings of non-uniform
%           thickness, rr(ddy/dx²) = (ddy/dt²), and the case
%           y = xΓ:(β/x − t) + xΔ:(β/x + t). Runs on to v. 742.
%   v. 742  (ink 369) end of the Euler sentence (isochronous vibrations);
%           then « La Grange Pag. 259 du même volume »: the initial curve
%           and y = α sin πx + β sin 2πx + &c., « la même géométriquement »,
%           the initial curve as an asymptote of the generating curve;
%           y = α sin πx/a + β sin 2πx/b + γ sin 3πx/c (denominators read
%           a, b, c as written); « Pag. 275 », Δr = Δ(a+iy) expanded,
%           « [cequi est bien près de la théorie des fonctions.] ».
%   v. 743  (verso of 369) a new extract under a title: « Mémoire de La
%           Grange sur l'intégration de quelques équations differentielles
%           dont les indéterminées sont séparées, mais dont chaque membre
%           en particulier n'est point intégrable — T. 4 »: equation (A)
%           dx/√(1−x²) = dy/√(1−y²), integral (B) x = y√(1−a²) + a√(1−y²).
%   v. 744  (ink 370) the same integral by imaginary logarithms, then by
%           a purely algebraic route (multiplying across and integrating
%           by parts).
%   v. 745  (verso of 370) the method applied to dx/√(α+βx+γx²) =
%           dy/√(α+βy+γy²) (C), (D), and the algebraic integral
%           (x + β/2γ)√(α+βy+γy²) = (y + β/2γ)√(α+βx+γx²) + C; the
%           quartic case announced; the extract closes with a bracketed
%           note that the rest of the memoir (« environ 20 pages ») is
%           copied almost whole « depuis la page 77 jusqu'a la p. 91 du
%           2me V. du C. I. de Cousin » (reading of the author's name
%           doubtful).
%   v. 746  (ink 371) « Condorcet Pag. 4 du 4me V. de Turin » (both 4s
%           doubtful): Z an exact differential, δdB = δZ, ∫δZ = δB, the
%           integration by parts of δZ; then « Pag 168 et suivantes
%           Memoire de La Grange sur la methode des variations »: ψ, Π,
%           ∫ψ + Π = a constant, and a bracketed comparison of Condorcet
%           and Lagrange ending « il faudroit avoir un mémoire précédent
%           sous les yeux ».
%   v. 747  (verso of 371) a bracketed note that pp. 188–243 hold two
%           memoirs of Lagrange « qu'il a analysés » in « la section de la
%           mechanique » on the motion of a body attracted to one or
%           several centres; « Pag. 210 », two fixed centres, one sent to
%           infinity; « Pag. 271 » of the Mécanique; « 2ème Memoire Pag.
%           218 »: d²x/2dt² + M = 0, d²y/2dt² + N = 0, the multipliers
%           m dx + n dy, μ dx + ν dy, and the form of the first integral.
%           Runs on to v. 748.
%   v. 748  (ink 372) the conditions on A, B, C: A = a + by + cy²,
%           C = g + hx + cx², B = k − bx − hy − 2cxy, with bracketed
%           justifications, one of them heavily rewritten between the
%           lines; the condition d(2AM+BN)/dy = d(BM+2CN)/dx. Runs on to
%           v. 749.
%   v. 749  (verso of 372; stamp « BIBLIOTHEQUE ROYALE ») end of the same
%           memoir: M and N in terms of D = 4AC − B² and an arbitrary V,
%           the integral (A dx² + B dx dy + C dy²)/2dt² + D²V = const.
%           Then, under a heavy rule, NUMBER THEORY: « [Dans le mémoire
%           cité (P. 37) dela Théorie des nombres on trouve (P. 46 T. 4
%           de Turin) » — the Pell equation x² − ay² = 1 after Lagrange:
%           the lemma on the product (x² − ay²)(x'² − ay'²), the powers
%           (p ± q√a)^m giving infinitely many solutions from one, the
%           closed forms x = [(p+q√a)^m + (p−q√a)^m]/2, y = […]/2√a, and
%           « La Grange demontre ensuite P. 65 » that the least solution
%           generates all. This is Lagrange's classical theory, copied;
%           it is not Fermat's Last Theorem and not Manuscript C. The
%           written part of the volume ends here, mid-page.
%   v. 750  a blank leaf mounted on its guard, in front of the back
%           endpapers: no text; faint traces of writing show through from
%           the other side (not photographed). Top right, in ink, a
%           number read as « 373 », struck through with two horizontal
%           strokes. No « Volume de N feuillets » certificate or other
%           library note is visible on this view. Not transcribed as
%           text; recorded in the opening note.
%
% Numbers on the leaves. The top-right ink series runs 369 (v. 742), 370
% (v. 744), 371 (v. 746), 372 (v. 748), continuing 368 on v. 740, one per
% recto, none on the versos; then a struck « 373 » on the blank leaf of
% v. 750. The last unstruck number, 372, is exactly the leaf count the
% catalogue gives (« Papier. 372 feuillets »): the series ends where the
% library's count ends, and the one number beyond it is cancelled on a
% blank leaf. This is evidence for the volume's open foliation question
% (see hand.md), not a settlement: the figures are still penned, and no
% \folio{} is written, by the house rule. No page or paragraph numbers
% of the copyist's. Other numbers are page references to printed
% volumes (Turin vol. 3–4 pp. 259, 275, 168, 188–243, 210, 218; « T. 4 »;
% « P. 37 », « P. 46 », « P. 60 », « P. 65 »; « la méchanique (Pag.
% 271) »; Cousin p. 77–91), transcribed where they stand.
%
% Notation. Euler's functional colon kept (Γ:(x+ct)); the coefficient
% letters drawn as « u », a looped ℓ and « v » are set α, β, γ, as
% elsewhere in the volume; π drawn as two minims; ψ in the variations
% extract is a barred « v » (\uncertain at first use); Π drawn as two
% minims under a bar; δ a looped d. Underlined words are set \emph{}.
% The copyist's square brackets are kept as square brackets. Spelling as
% written (« sauroit », « tems », « fontion », « quarré », « fesant »,
% « étants », « multipicateurs », run-together « dela », « demême »,
% « aulieu », « engénéral », « cequi »).
%
% \ill{} in this batch: five, all inside a deletion — a struck word
% before « ainsi » (v. 743); a word under a blot after « indiquée » and a
% short struck word after « disparoissent » (v. 748); a struck word
% before « naturels » and one before « q » (v. 749).
%
% Across the edges. v. 741 continues the recto of v. 740 (ink 368, Euler's
% « Eclaircissemens », batch 37). Nothing follows v. 750: last view.
%
\documentclass[11pt,a4paper]{article}
% The preamble shared by every transcription of the mathematicians' manuscripts in Gallica.
%
% It rests on one idea: **uncertainty compounds, it does not resolve.** A
% transcription that smooths over a doubtful word has destroyed the only thing
% it was made for — a reader must be able to tell, at every point, what was on
% the page and what was supplied. The apparatus macros below are therefore the
% critical apparatus, not decoration.
%
% The same commands are recognised by `scripts/render.mjs`, which produces the
% site's reading view, and by `scripts/tei.mjs`. A macro added here must be
% added there too, or rendering fails loudly — which is the wanted behaviour.
%
% Comments here are English, like the rest of the repository. The documents
% this preamble sets are French, and that is deliberate: see the skills.

\ProvidesPackage{germain}[2026/09/15 Transcriptions of mathematicians' manuscripts in Gallica]

\usepackage{amsmath,amssymb,amsthm}
\usepackage{mathrsfs}
% Kept from the parent preamble so the renderer's diagram support stays
% available; these manuscripts draw no commutative diagrams, and nothing here uses it.
\usepackage{tikz-cd}
\usepackage[english,french]{babel}
\usepackage{fontspec}

% --- Greek ------------------------------------------------------------------
% The text face stays fontspec's default, Latin Modern, which has no Greek:
% XeTeX drops every glyph it lacks with a line in the log and nothing on the
% page, and the Greek words the manuscripts quote vanished from the PDFs. So
% the Greek and Greek Extended blocks get a XeTeX character class of their own,
% and entering or leaving a run of it switches the family to CMU Serif — the
% same Computer Modern design, with polytonic Greek — and back. Only the
% family changes: a Greek word in \textit or \textbf keeps its shape and series.
% The transitions to and from every other class carry the tokens that class
% has towards an ordinary letter, so French spacing before « ; » or inside
% « » still holds after a Greek word. No grouping, so no brace or \endgroup
% can come out unbalanced; maths is left alone, since XeTeX inserts nothing
% there. Kept identical in `exercices.sty`.
\makeatletter
\newfontfamily\germainGreekfont{cmun}[
  Extension=.otf, NFSSFamily=germaingreek,
  UprightFont=*rm, ItalicFont=*ti, SlantedFont=*sl,
  BoldFont=*bx, BoldItalicFont=*bi, BoldSlantedFont=*bl]
\newXeTeXintercharclass\germain@greek
\def\germain@greekrange#1#2{%
  \@tempcnta=#1\relax
  \loop
    \XeTeXcharclass\@tempcnta=\germain@greek
    \ifnum\@tempcnta<#2\advance\@tempcnta\@ne\repeat}
\germain@greekrange{"0370}{"03FF}
\germain@greekrange{"1F00}{"1FFF}
\def\germain@greekprev{\rmdefault}
\def\germain@greekin{%
  \edef\germain@greekprev{\f@family}\fontfamily{germaingreek}\selectfont}
\def\germain@greekout{\fontfamily{\germain@greekprev}\selectfont}
\def\germain@greekjoin{%
  \XeTeXinterchartokenstate=\@ne
  \@tempcnta=\z@
  \loop
    \ifnum\@tempcnta=\germain@greek\else
      \XeTeXinterchartoks\germain@greek\@tempcnta=\expandafter{%
        \expandafter\germain@greekout\the\XeTeXinterchartoks\z@\@tempcnta}%
      \XeTeXinterchartoks\@tempcnta\germain@greek=\expandafter{%
        \the\XeTeXinterchartoks\@tempcnta\z@\germain@greekin}%
    \fi
    \ifnum\@tempcnta<4095\advance\@tempcnta\@ne\repeat}
% babel-french sets its own classes, and saves and restores the interchar
% state, each time a language is selected: join again after it does.
\addto\extrasfrench{\germain@greekjoin}
\addto\extrasenglish{\germain@greekjoin}
\AtBeginDocument{\germain@greekjoin}
\makeatother

\usepackage{geometry}
\geometry{a4paper,margin=25mm,marginparwidth=28mm}
\usepackage{marginnote}
\usepackage{xcolor}
\usepackage[normalem]{ulem}
\setlength{\emergencystretch}{3em}
\usepackage{hyperref}
\hypersetup{colorlinks=true,linkcolor=black,urlcolor=black,pdfborder={0 0 0}}

% --- Batch metadata ---------------------------------------------------------
% Taken as they stand by the rendering script, which shows them at the head of
% the reading view. They fix what the file claims to cover. The macro names are
% the parent project's, so that the scripts stay shared; \folder here means the
% volume's slug — `fr-9115` — and \pages counts Gallica views.

\newcommand{\folder}[1]{\gdef\grothFolder{#1}}
\newcommand{\batch}[1]{\gdef\grothBatch{#1}}
\newcommand{\pages}[2]{\gdef\grothFirst{#1}\gdef\grothLast{#2}}
\newcommand{\foldertitle}[1]{\gdef\grothFoldertitle{#1}}

% The holder's own shelfmark, printed as they write it: « Français 9115 ».
\newcommand{\shelfmark}[1]{\gdef\grothShelfmark{#1}}

% The dating, copied verbatim from the holder's catalogue — never inferred
% here. « XIXe siècle » is the BnF's; a pass that dates a leaf from its content
% says so in a \note{}, as its own inference.
\newcommand{\dating}[1]{\gdef\grothDating{#1}}

% The status notice, printed in the title block and repeated as a diagonal
% watermark on every page of the PDF. The manuscripts are in the public
% domain; what this line declares is not a right but a quality — an
% unverified first pass by a machine. The skills write the line; a file
% without it gets no watermark, so leaving it out is visible in the output.
\newcommand{\watermark}[1]{\gdef\grothWatermark{#1}}

\gdef\grothFolder{?}\gdef\grothBatch{}\gdef\grothFirst{?}\gdef\grothLast{?}%
\gdef\grothFoldertitle{}\gdef\grothShelfmark{}\gdef\grothDating{}\gdef\grothWatermark{}

% --- The résumé -------------------------------------------------------------
\newenvironment{resume}
  {\small\setlength{\parskip}{0.45em}}
  {\par\medskip\hrule\medskip}

% --- Keywords ---------------------------------------------------------------
% In English, closing the modernised reading's résumé: the modern vocabulary
% under which to search for what the leaves construct. The manifest extracts
% them to tag the volume — this line is the single source of the tags.
\newcommand{\keywords}[1]{\par\smallskip\noindent
  {\footnotesize\textsc{Keywords} --- \textit{#1}}\par}

% --- The critical apparatus -------------------------------------------------

% The Gallica view number, in the margin. This is the anchor that lets a
% disagreement over a formula be settled by looking at *one* image rather than
% a volume of 750. The site uses it to turn the facsimile as the transcription
% is scrolled. A view is one scanned image: usually one side of a leaf.
\newcommand{\page}[1]{%
  \marginnote{\footnotesize\color{gray!70}\textsf{v.\,#1}}%
  \label{page:#1}\ignorespaces}

% The BnF's pencilled foliation, where the leaf carries it and a pass can read
% it: « 198r ». This is what the literature cites — « ff. 198r–208v » — and
% the one line that turns a citation into a view. Recorded, never inferred:
% a folio a pass cannot read is not written.
\newcommand{\folio}[1]{%
  \marginnote{\footnotesize\color{gray!50}\textsf{f.\,#1}}\ignorespaces}

% The modernised reading's coarser counterpart to \page{}: which views a
% section is drawn from.
\newcommand{\pagerange}[2]{%
  \marginnote{\footnotesize\color{gray!70}\textsf{v.\,#1--#2}}%
  \ignorespaces}

% Illegible: never guessed.
\newcommand{\ill}{\textcolor{red!60!black}{[\,\ldots\,]}}

% A doubtful reading: the word is given, and flagged as uncertain.
\newcommand{\uncertain}[1]{\uline{#1}}

% An editorial addition — a missing letter, an implied word.
\newcommand{\add}[1]{\textcolor{blue!50!black}{[#1]}}

% Struck out by the writer. What was crossed out is part of the document.
\newcommand{\struck}[1]{\sout{#1}}

% The transcriber's note, on the state of the page or the mathematics.
\newcommand{\note}[1]{\par\smallskip{\small\color{gray!60!black}\textit{[#1]}}\par\smallskip}

% A marginal note of hers, distinct from the body of the text.
\newcommand{\marginal}[1]{\par\smallskip\noindent\hspace{1em}%
  {\small\color{gray!60!black}#1}\par\smallskip}

% --- Title ------------------------------------------------------------------
% The title block is French, like the documents themselves.

\AddToHook{shipout/background}{%
  \ifx\grothWatermark\empty\else
    \begin{tikzpicture}[remember picture, overlay]
      \node[rotate=52, scale=3.4, align=center, text=gray!18,
            font=\sffamily\bfseries]
        at (current page.center) {\grothWatermark};
    \end{tikzpicture}%
  \fi}

\AtBeginDocument{%
  \begin{center}
    {\large\bfseries Manuscrits de mathématiciens (Gallica) ---
      \ifx\grothShelfmark\empty\grothFolder\else\grothShelfmark\fi}\\[2pt]
    {\itshape\grothFoldertitle}\\[4pt]
    {\small vues \grothFirst--\grothLast
      \ifx\grothBatch\empty\else\ (lot \grothBatch)\fi}\\[2pt]
    \ifx\grothDating\empty\else
      {\small\color{gray!60!black}Datation du catalogue : \grothDating}\\[2pt]
    \fi
    \ifx\grothWatermark\empty\else
      {\small\bfseries\color{red!45!gray}\grothWatermark}\\[2pt]
    \fi
    {\footnotesize\color{gray!60!black}Transcription établie d'après les images IIIF de Gallica.
     Source gallica.bnf.fr / Bibliothèque nationale de France.\\
     Les lectures incertaines sont soulignées, les illisibles notées [\,\ldots\,],
     les ajouts entre crochets ; \textsf{v.} numérote les vues de Gallica,
     \textsf{f.} la foliotation au crayon de la BnF.}
  \end{center}
  \vspace{1em}}

\endinput


\folder{fr-9115}
\shelfmark{Français 9115}
\batch{38}
\pages{741}{750}
\dating{1801-1900}
\watermark{Lecture automatique\\non vérifiée}
\foldertitle{Papiers de Sophie GERMAIN. — Recueil de dissertations et problèmes mathématiques et physiques.}

\begin{document}

\note{Numérotation des feuillets. En haut à droite, à l'encre, d'une écriture plus ronde, la série qui court dans tout le volume : 369 (vue 742), 370 (744), 371 (746), 372 (748), sur le recto seulement, à la suite de 368 (vue 740, hors de ce lot) ; les versos (vues 741, 743, 745, 747, 749) ne portent pas de numéro. Sur la vue 750, feuillet blanc qui clôt le volume, en haut à droite, un nombre à l'encre lu « 373 », barré de deux traits horizontaux. Le dernier nombre non barré, 372, est le nombre de feuillets que donne la notice du volume (« 372 feuillets ») ; la série occupe la place de la foliotation de la Bibliothèque, mais elle est tracée à la plume : aucun « f. » n'est donc porté. Aucune mention du type « Volume de \ldots{} feuillets » n'est visible sur la dernière vue. Tous ces nombres ont été lus sur la vue ; aucun n'est déduit. La vue 750 n'est pas transcrite : feuillet blanc, où transparaissent quelques traces d'écriture de l'autre face, non photographiée.}

\page{741}

\note{Verso du feuillet de la vue 740 (encre 368, en haut à droite), dont le recto, hors de ce lot, porte le titre « 3 V. des Memoires de Turin --- Eclaircissemens sur le mouvement des cordes vibrantes Par Euler ». La page commence un nouvel alinéa, à la suite de ce recto. Aucun numéro écrit de ce côté.}

On ne sauroit douter non plus qu'après avoir imprimé a la corde une telle figure
discontinue ou irréductible a aucune équation analytique, la corde étant subitement
relachée, soit qu'on lui imprime encore quelque mouvement ou non, n'en recoive
un certain mouvement de vibration.

Mais si la théorie nous conduit à une solution si générale, qu'elle
s'etend aussi bien à toutes les figures discontinues que continues il faudra
avouer, que cette recherche ouvre une nouvelle carrière dans l'analyse.
En effet j'ai remarqué a cette occasion, que la partie de l'analyse
des infinis, a laquelle cette question appartient, renferme essentiellement
ce caractère, qu'elle recoit \uncertain{dele} fonctions absolument arbitraires, pendant
que de telles fonctions sont entièrement bannies de l'analyse ordinaire
qu'on a cultivée jusqu'ici, qui roule principalement sur des fonctions
d'une seule variable. Mais l'analyse dont nous avons besoin ici,
s'occupe des fonctions de deux ou plusieurs variables: ou cela arrive
de bien remarquable, que chaque intégration introduit dans le
calcul une fonction absolument arbitraire.
\note{« qu'elle recoit \uncertain{dele} fonctions » : le mot, en un seul trait, se lit « dele » ; on attendrait « des ». Après « cultivée », un petit signe en forme de croix, peut-être un renvoi sans note correspondante sur la page. « ou cela arrive » : ainsi, sans accent.}

[c'est cequ montre l'équation $y = \Gamma{:}(x+ct) + \Delta{:}(x-ct)$ ou $\Gamma{:}(x+ct)$ marque
une fonction quelconque de la quantité $x+ct$, $\Delta{:}(x-ct)$ une fontion aussi
quelconque de $(x-ct)$, qui l'intégrale de $cc\left(\frac{ddy}{dx^{2}}\right) = \left(\frac{ddy}{dt^{2}}\right)$ equation aux
cordes vibrates]
\note{Les crochets sont sur la feuille : ce passage entre crochets paraît une explication ajoutée à la traduction, non le texte d'Euler ; c'est une inférence. « cequ », « fontion », « qui l'intégrale » (pour « qui est l'intégrale ») et « vibrates » : ainsi écrits. Le deux-points après $\Gamma$ et $\Delta$ est la notation des fonctions, conservée.}

Lorsque l'on veut considerer des cordes qui ne soient pas uniformement
epaisses on a aulieu de $cc\left(\frac{ddy}{dx^{2}}\right) = \left(\frac{ddy}{dt^{2}}\right)$ dans laquelle $c$ est une
constante; $rr\left(\frac{ddy}{dx^{2}}\right) = \left(\frac{ddy}{dt^{2}}\right)$ et \uncertain{$r$} represente une fontion de $x$ donnée
par la grosseur et la tension de la corde l'intégrale de cette
équation est engénéral de la forme
\[ y = P\Gamma{:}\left(\int u\,dx + t\right) + Q\Gamma'{:}\left(\int u\,dx + t\right) + R\Gamma''{:}\left(\int u\,dx + t\right) + \&c. \]
ou $P$, $Q$, $R$ \&c. demême que $u$ sont de certaines fonctions dela seule
variable $x$. Ensuite puisque la quantité $t$ se peut prendre tant
aussi bien negativement que positivement, pour compléter l'intégrale
on y doit ensuite ajouter encore ces termes
\[ P\Delta{:}\left(\int u\,dx - t\right) + Q\Delta'{:}\left(\int u\,dx - t\right) + R\Delta''{:}\int\left(u\,dx - t\right) + \&c. \]
qui n'exigent donc point une recherche particulière
\note{« et \uncertain{$r$} represente » : entre « et » et « represente » un $r$ se lit mal, comme redoublant l'initiale du verbe. « tant aussi bien » : ainsi. Dans le dernier terme, le signe $\int$ est écrit devant la parenthèse, et non dedans comme dans les deux premiers ; conservé. Les signes $-$ de la seconde série sont un trait suivi d'un point, qu'on lit « $-$ » par l'opposition annoncée entre $t$ positif et négatif.}

Dans le cas ou on a simplement
\[ y = x\Gamma{:}\left(\frac{\beta}{x} - t\right) + x\Delta{:}\left(\frac{\beta}{x} + t\right) \quad \text{et ou parconséquent} \]
\[ P = x, \quad \text{et} \quad \int u\,dx = -\frac{\beta}{x} \]
les cordes ont cela de commun avec les cordes uniformement epaisses,
que de quelque manière qu'on les ébranle, elles produisent toujours des
vibrations regulières. Le fondement de cette proprieté est, qu'en
quelqu'état que la corde soit mise au commencement elle y doit
necessairement revenir après un certain tems; or dans les autres
espèces de cordes ce rétablissement dans le premier état n'arrive
\note{La phrase se poursuit à la vue 742.}

\page{742}

\note{En haut à droite, à l'encre, « 369 ». Le feuillet est monté bas sur sa feuille de garde ; au-dessus de son bord, sur la tranche d'un autre feuillet qui dépasse, se lit « 370 » : c'est le numéro du feuillet suivant (vue 744), non de celui-ci. La phrase de la vue 741 continue en tête de page.}

peut-être jamais, et quand il arrive ce n'est que sous de certaines
circonstances sous lequel le son ne sauroit être regulier ou formé
par des \struck{vi} vibrations \emph{isochrones}.
\note{« sous lequel » : ainsi, au singulier. « isochrones » est souligné. Le passage tiré d'Euler s'achève sur cette phrase ; la suite est donnée sous le nom de Lagrange.}

La Grange \emph{Pag. 259} du même volume

On pourroit douter s'il ne faut pas que la courbe initiale de la corde
soit comprise aussi dans l'équation $y = \alpha\,\text{sin.}\,\pi x + \beta\,\text{sin.}\,2\pi x + \&c.$
Il est certain que si on veut que la courbe génératrice soit la même
\emph{géométriquement} que la courbe initiale, il faut que celle-ci soit
renfermé dans l'équation $y = \alpha\,\text{sin}\,\pi x + \&c.$ Je dis la même \emph{géométriquement}
car il suffit que la difference de ces deux courbes soit moindre
qu'aucune grandeur donné pour qu'elles puissent être prises pour
les mêmes. Or il est clair que, quelque soit la courbe initiale
on peut toujours faire passer, par une infinité de points infiniment
proches de cette courbe, une autre courbe de la forme
$y = \alpha\,\text{sin}\,\pi x + \&c.$ de manière que la difference entre ces deux
courbes soit aussi petite que l'on voudra, quoique cette difference
ne puisse devenir absolument nulle que dans le cas ou la courbe
initiale sera aussi dela même forme; dans tous les autres cas cette
courbe initiale ne sera qu'une espèce d'assymptôte dont la courbe
génératrice pourra s'approcher a l'infini, sans qu'elle puisse jamais
coincider entièrement.
\note{« Pag. 259 » est souligné ; « du même volume » renvoie, semble-t-il, au volume des Mémoires de Turin nommé en tête de la vue 740. Le $\pi$ est écrit comme deux jambages, « $11$ » ; le premier coefficient comme un « a » ou « u », le second comme un $\beta$ fermé. « donné », « renfermé » : ainsi, sans accord. Une tache d'encre sur « nulle ».}

On demontre dans le même memoire que l'équation dela
figure initiale de la corde lorsqu'elle en a une, ne peut être
que de la forme
\[ y = \alpha\,\text{sin.}\,\frac{\pi x}{a} + \beta\,\text{sin.}\,\frac{2\pi x}{b} + \gamma\,\text{sin.}\,\frac{3\pi x}{c} + \&c. \]
\note{Les dénominateurs se lisent $a$, $b$, $c$, et non trois fois la même longueur ; transcrits tels. Le troisième coefficient est tracé comme un « $\vartheta$ », qui est son $\gamma$ ailleurs dans le volume.}

A la \emph{Pag. 275} du même mémoire on y trouve l'équation $r = a + iy$
$i$ étant un coëfficient très petit et $y$ une nouvelle variable et par conséquent
\[ \Delta r = \Delta(a + iy) = \left(\text{en supposant } \frac{d.\Delta r}{dr} = \Delta' r,\ \frac{d.\Delta' r}{dr} = \Delta'' r\ \&c.\right) \]
\[ = \Delta a + i\Delta' a\,y + i^{2}\frac{\Delta'' a}{2}\,y^{2} + i^{3}\frac{\Delta''' a}{2.3}\,y^{3} + \&c. \]

[cequi est bien près de la théorie des fonctions.]
\note{« Pag. 275 » : le mot est souligné. Le crochet final est sur la feuille, comme à la vue 741 : remarque du copiste sur la page de Lagrange, rapprochée de la théorie des fonctions analytiques ; l'attribution est une inférence. Le bas de la page est blanc.}

\page{743}

\note{Verso du feuillet 369 (vue 742) ; aucun numéro écrit de ce côté : ce qui paraît en haut à gauche est le « 369 » du recto lu par transparence, à rebours. Un nouveau morceau commence ici, sous un titre.}

Mémoire de La Grange sur l'intégration de quelques équations
differentielles dont les indéterminées sont séparées, mais dont
chaque membre en particulier n'est point intégrable

T. 4
\note{« T. 4 », sous le titre, au-dessus d'un trait : sans doute le tome du recueil où se trouve le mémoire ; c'est une inférence.}

La séparation des indéterminées est regardée avec raison comme un des
meilleurs moyen que les géometres ayent imaginés pour intégrer les
équations differentielles du premier ordre. En effet il est clair
que quand on a séparé les indéterminées dans une équation, on
peut alors regarder chacun des membres comme une differentielle
particulière qui ne contient qu'une variable; de sorte qu'il n'y
a plus qu'a prendre séparement l'intégrale de l'un et de
l'autre membre, en y ajoutant une constante arbitraire.
Delà il semble qu'on pourroit conclure que lorsque les deux
membres de l'équation \struck{\ill{}} ainsi separée ne sont point integrables
l'équation elle-même ne doit pas l'etre non plus; c'est ce
qui est vrai en effet dans la plus part des équations
differentielles; mais il se trouve néanmoins des cas, ou cette
conclusion seroit fausse, et qui vont faire le sujet de ce
\emph{mémoire}.
\note{« un des meilleurs moyen » : ainsi, au singulier. Devant « ainsi », un mot bref barré d'un trait, qu'on ne lit pas (quelque chose comme « iute »). « mémoire », à la fin de l'alinéa, est souligné.}

Pour commencer par les cas les plus simples nous prendrons
l'équation \quad $\frac{dx}{\sqrt{(1-x^{2})}} = \frac{dy}{\sqrt{(1-y^{2})}} \ldots\ldots (\text{A})$
dans laquelle tout est separé comme l'on voit. Il est d'abord
evident que les deux membres de cette équation ne sont point intégrables
au moins algébriquement; cependant on sait que l'équation en elle-même
admet une intégrale algébrique. En effet $\frac{dx}{\sqrt{(1-x^{2})}}$ est la
differentielle de l'arc dont le le sinus est $x$, demême que $\frac{dy}{\sqrt{(1-y^{2})}}$
est la differentielle de l'arc dont le sinus est $y$ on aura en
prenant les arcs aulieu de leurs differentielles et ajoutant une
constante quelconque $C$ et supposant $C = \text{arc. dont le sinus est } a$
\[ \text{arc sin.}\,x = \text{arc sin.}\,y + \text{arc sin.}\,a \]
c'est-adire que l'arc qui repond au sinus $x$ doit être égal à la
somme des arcs qui repondent aux sinus $y$ et $a$. de sorte que
l'on aura par les théorèmes connus
\[ x = y\sqrt{(1-a^{2})} + a\sqrt{(1-y^{2})} \ldots\ldots (\text{B}) \]
[en effet on a engénéral $\text{sin.}(m+n) = \text{sin.}\,n\,\text{cos.}\,m + \text{sin.}\,m\,\text{cos.}\,n$ et il est visible
que $\sqrt{(1-a^{2})}$ et $\sqrt{(1-y^{2})}$ representent les cosinus des arcs dont les sinus sont $a$ et $y$]
C'est l'intégrale de la proposée, dans laquelle $a$ est la constante
arbitraire.
\note{« dont le le sinus » : le « le » est écrit deux fois, sans rature. Dans (B), le premier $a$ est écrit comme un $\alpha$. Le crochet est sur la feuille, comme aux vues 741 et 742. Le texte s'arrête au bas du feuillet et reprend à la vue 744.}

\page{744}

\note{En haut à droite, à l'encre, « 370 » : le numéro qui dépassait au-dessus du feuillet de la vue 742. Le texte suit, sans rupture, celui de la vue 743 (mémoire de Lagrange, équation (A)).}

J'avoue qu'on peut trouver cette intégrale sans le secours des théorèmes sur les
sinus, en intégrant chaque membre de l'équation (A) par les logarithmes
imaginaires, et passant ensuite des logarithmes aux nombres de cette manière
on aura [en mettant $\frac{dx}{\sqrt{(1-x^{2})}} = \frac{dy}{\sqrt{(1-y^{2})}}$ sous la forme
\[ \frac{dx\left[\sqrt{(1-x^{2})} + x\sqrt{-1}\right]}{\sqrt{(1-x^{2})}\left[\sqrt{(1-x^{2})} + x\sqrt{-1}\right]} = \frac{dy\left[\sqrt{(1-y^{2})} + y\sqrt{-1}\right]}{\sqrt{(1-y^{2})}\left[\sqrt{(1-y^{2})} + y\sqrt{-1}\right]} \]
\marginal{multipliant les deux membres par $\sqrt{-1}$}
\[ \frac{dx\sqrt{-1}}{\sqrt{(1-x^{2})} + x\sqrt{-1}} - \frac{x\,dx}{\sqrt{(1-x^{2})}\left[\sqrt{(1-x^{2})} + x\sqrt{-1}\right]} \]
\[ = \frac{dy\sqrt{-1}}{\sqrt{(1-y^{2})} + y\sqrt{-1}} - \frac{y\,dy}{\sqrt{(1-y^{2})}\left[\sqrt{(1-y^{2})} + y\sqrt{-1}\right]} \]
ou \quad $\left[dx\sqrt{-1} - \frac{x\,dx}{\sqrt{(1-x^{2})}}\right] : \sqrt{(1-x^{2})} + x\sqrt{-1} = \left[dy\sqrt{-1} - \frac{y\,dy}{\sqrt{(1-y^{2})}}\right] : \sqrt{(1-y^{2})} + y\sqrt{-1}$
\[ d.L\left[x\sqrt{-1} + \sqrt{(1-x^{2})}\right] = d.L\left[y\sqrt{-1} + \sqrt{(1-y^{2})}\right] \]
puis intégrant en observant que la constante doit être telle que le premier
membre de l'équation lui devienne egal lorsque $y = 0$]
\note{« multipliant les deux membres par $\sqrt{-1}$ » est écrit dans la marge de gauche, en face de la deuxième équation. Le deux-points de la troisième ligne marque la division. Le premier crochet de $d.L[\ldots$ se ferme sur la feuille par une parenthèse.}

\[ l\left[x\sqrt{-1} + \sqrt{(1-x^{2})}\right] = l\left[y\sqrt{-1} + \sqrt{(1-y^{2})}\right] + l\left[a\sqrt{-1} + \sqrt{(1-a^{2})}\right] \]
d'ou l'on tire
\[ x\sqrt{-1} + \sqrt{(1-x^{2})} = \left[y\sqrt{-1} + \sqrt{(1-y^{2})}\right]\left[a\sqrt{-1} + \sqrt{(1-a^{2})}\right] \]
\[ = \left[y\sqrt{(1-a^{2})} + a\sqrt{(1-y^{2})}\right]\sqrt{-1} + \sqrt{(1-y^{2})}\sqrt{(1-a^{2})} - ay \]
et comparant la partie imaginaire du premier membre à la partie
imaginaire du second, et la partie réelle avec la partie réelle; on aura
comme cidessus
\[ x = y\sqrt{(1-a^{2})} + a\sqrt{(1-y^{2})} \]
ou bien encore ce qui revient au même dans le fond.
\[ \sqrt{(1-x^{2})} = \sqrt{(1-a^{2})}\sqrt{(1-y^{2})} - ay \]
[cette dernière équation resulte de la comparaison des parties réelles \add{mais on}
peut la tirer \add{aussi} dela première; en effet si on la quarre que l'on ajoute \add{ensuite} l'unité
a chacun de ses membres \add{pris avec un signe contraire} et que l'on prenne \struck{ensuite} la racine quarrée
on aura \quad $\sqrt{(1-x^{2})} = \sqrt{\left(1 - y^{2}(1-a^{2}) - 2ya\sqrt{(1-y^{2})}\sqrt{(1-a^{2})} - a^{2}(1-y^{2})\right)}$
\[ = \sqrt{\left((1-a^{2})(1-y^{2}) - 2ya\sqrt{(1-y^{2})}\sqrt{(1-a^{2})} + a^{2}y^{2}\right)} \]
\[ = \sqrt{(1-a^{2})}\sqrt{(1-y^{2})} - ay\ ] \]
\note{« mais on », d'une encre plus appuyée, est écrit en bout de ligne après « réelles », peut-être sur un mot effacé ; « aussi », « ensuite » (le premier) et « pris avec un signe contraire » sont écrits au-dessus de la ligne, en plus petit ; le second « ensuite », dans la ligne, est barré. Le crochet ouvert devant « cette dernière » se ferme après la dernière égalité.}

Mais si d'un côté cette méthode est un peu plus directe que la
précédente, de l'autre elle a aussi l'inconvenient de dependre des
quantités transcendantes; en effet puisque l'intégrale de l'équation
est purement algébrique n'est-il pas naturel depenser qu'il y ait
une voie \add{aussi} purement algébrique pour y parvenir?
Qu'on multiplie les deux membres de l'équation (A) en croix, on aura
\[ dx\sqrt{(1-y^{2})} = dy\sqrt{(1-x^{2})} \quad \text{et intégrant par parties} \]
\[ x\sqrt{(1-y^{2})} + \int\frac{xy\,dy}{\sqrt{(1-y^{2})}} = y\sqrt{(1-x^{2})} + \int\frac{yx\,dx}{\sqrt{(1-x^{2})}} + C \]
or l'équation (A) étant multipliée par $yx$ et ensuite integrée donne
$\int\frac{xy\,dy}{\sqrt{(1-y^{2})}} = \int\frac{xy\,dx}{\sqrt{(1-x^{2})}}$ donc l'équation précédente deviendra $x\sqrt{(1-y^{2})} = y\sqrt{(1-x^{2})} + C$
équation algébrique qui en déterminant convenablement la constante revient au
même que celle trouvée plus haut.
\note{« aussi », au-dessus de « purement », est ajouté entre les lignes. Dans le second membre de la deuxième équation, un premier $y$ est écrit et chargé, puis repris : on lit $yx\,dx$. « multipliée » : le « m » est récrit sur un autre trait. La page s'arrête au bas du feuillet sur « plus haut. » et la vue 745 continue.}

\page{745}

\note{Verso du feuillet 370 (vue 744) ; aucun numéro écrit de ce côté, le « 370 » du recto se lit à rebours par transparence en haut à gauche. La suite immédiate de la vue 744. Les coefficients sont tracés comme ailleurs dans le volume : $\alpha$ comme un « u », $\beta$ comme un « $\ell$ » bouclé, $\gamma$ comme un « v ».}

On pourroit aussi appliquer la même méthode a l'intégration de
l'équation générale \quad $\frac{dx}{\sqrt{(\alpha + \beta x + \gamma x^{2})}} = \frac{dy}{\sqrt{(\alpha + \beta y + \gamma y^{2})}} \ldots\ldots \text{C}$
car multipliant d'abord en croix, et prenant ensuite l'intégrale de
chaque membre par parties on a
\[ x\sqrt{(\alpha + \beta y + \gamma y^{2})} - \int\frac{(\beta + 2\gamma y)\,x\,dy}{2\sqrt{(\alpha + \beta y + \gamma y^{2})}} \]
\[ = y\sqrt{(\alpha + \beta x + \gamma x^{2})} - \int\frac{(\beta + 2\gamma x)\,y\,dx}{2\sqrt{(\alpha + \beta x + \gamma x^{2})}} + C \ldots\ldots \text{D} \]
Or l'équation C étant multipliée par $xy$ et ensuite
intégrée donne
\[ \int\frac{xy\,dx}{\sqrt{(\alpha + \beta x + \gamma x^{2})}} = \int\frac{xy\,dy}{\sqrt{(\alpha + \beta y + \gamma y^{2})}} \]
deplus on a par la même équation
\[ \int\frac{y\,dx}{\sqrt{(\alpha + \beta x + \gamma x^{2})}} = \int\frac{y\,dy}{\sqrt{(\alpha + \beta y + \gamma y^{2})}} \]
\[ = \frac{1}{\gamma}\sqrt{(\alpha + \beta y + \gamma y^{2})} - \frac{\beta}{2\gamma}\int\frac{dy}{\sqrt{(\alpha + \beta y + \gamma y^{2})}} \]
et demême
\[ \int\frac{x\,dy}{\sqrt{(\alpha + \beta y + \gamma y^{2})}} = \int\frac{x\,dx}{\sqrt{(\alpha + \beta x + \gamma x^{2})}} \]
\[ = \frac{1}{\gamma}\sqrt{(\alpha + \beta x + \gamma x^{2})} - \frac{\beta}{2\gamma}\int\frac{dx}{\sqrt{(\alpha + \beta x + \gamma x^{2})}} \]
\[ = \frac{1}{\gamma}\sqrt{(\alpha + \beta x + \gamma x^{2})} - \frac{\beta}{2\gamma}\int\frac{dy}{\sqrt{(\alpha + \beta y + \gamma y^{2})}} \]
\note{Le signe « D » de la première équation est écrit en bout de ligne, sous le « $+ C$ ». Dans la troisième équation, le $y$ du second numérateur est récrit sur une autre lettre, sans doute un $x$.}

donc fesant ces substitutions dans l'équation (D) et effaçant
cequi se détruit on aura cette équation algébrique
\[ x\sqrt{(\alpha + \beta y + \gamma y^{2})} - \frac{\beta}{2\gamma}\sqrt{(\alpha + \beta x + \gamma x^{2})} = y\sqrt{(\alpha + \beta x + \gamma x^{2})} - \frac{\beta}{2\gamma}\sqrt{(\alpha + \beta y + \gamma y^{2})} + C \]
ou bien
\[ \left(x + \frac{\beta}{2\gamma}\right)\sqrt{(\alpha + \beta y + \gamma y^{2})} = \left(y + \frac{\beta}{2\gamma}\right)\sqrt{(\alpha + \beta x + \gamma x^{2})} + C \]
qui est l'intégrale de l'équation proposée
\note{« détruit » : le « t » est récrit au-dessus de la ligne.}

Voila donc comme on voit une methode bien simple
pour intégrer ces sortes d'équations, dont chaque membre en
particulier depend de la quadrature du cercle ou de
l'hyperbole; mais il y a encore d'autres equations plus
générales que les précédentes qui admettent aussi des intégrales
algébriques, quoique chacun de leurs membres ne soient en
aucune façon intégrable.
Ces équations sont comprises dans la formule suivante
\[ \frac{dx}{\sqrt{(\alpha + \beta x + \gamma x^{2} + \delta x^{3} + \varepsilon x^{4})}} = \frac{dy}{\sqrt{(\alpha + \beta y + \gamma y^{2} + \delta y^{3} + \varepsilon y^{4})}} \]
[le reste du memoire qui contient environ 20 pages est copié presqu'en entier
depuis la page 77 jusqu'a la p. 91 du 2\textsuperscript{me} V. \uncertain{du} C. I. de \uncertain{Cousin}]
\note{Le crochet final est sur la feuille : c'est une remarque de celui qui copie, qui renvoie pour la suite du mémoire à un autre ouvrage, dont le titre est abrégé « C. I. » (peut-être « Calcul intégral » ; c'est une inférence) ; le nom de l'auteur se lit mal. Le bas de la page est blanc ; le mémoire de Lagrange s'arrête ici.}

\page{746}

\note{En haut à droite, à l'encre, « 371 », au bout de la ligne de titre. Le feuillet est monté bas sur sa feuille de garde, comme celui de la vue 742. Nouveaux extraits, sans lien de phrase avec la vue 745.}

Condorcet \emph{Pag. \uncertain{4} du \uncertain{4}\textsuperscript{me} V. de Turin}
\note{Les deux chiffres ont la même forme bouclée, proche d'un « A », qui est le 4 du « T. 4 » de la vue 743 ; lecture donnée avec doute. La seconde moitié du titre est soulignée.}

Soit $Z$ une fonction de $x$, $y$, $z$ \&c. $dx$, $dy$, $dz$ \&c. $ddx$, $ddy$, $ddz$, \&c. \&c.
qu'elle soit une differentielle exacte d'une fonction d'un ordre moins
élevé, d'une unité, et que $B$ soit cette intégrale inconnue, il est
clair que differentiant d'une manière independante des differences qui se
trouvent dans $Z$ et dans $B$, j'ai $\delta dB = \delta Z$, \quad $\int\delta Z = \delta B$. le signe
d'intégration $\int$ se rapportant a la caractéristique $d$. Il suit dela
1.\textsuperscript{o}, que cherchant a intégrer $\delta Z$ par partie en regardant les $\delta x$, $\delta y$, $\delta z$ \&c.
comme des variables ordinaires, on aura necessairement égale a zéro
l'expression qui reste sous le signe $\int$ sans pouvoir être reduite.
[en effet on a lors $dZ$ a intégrer aulieu de $\delta Z$ et parconséquent l'intégrale
par le signe $\int$ n'est autre chose que $Z$, mais $Z$ a été supposée être
une differentielle exacte et une differentielle exacte ne doit contenir
aucune quantité sous le signe; donc la partie \struck{affect} affectée de
ce signe $\int$ doit être nulle]
2.\textsuperscript{o} que l'expression qui sera hors du signe sera identiquement la même
que $\delta B$ \ldots\ldots \&c.
\note{Le « \&c. » final de la première ligne est écrit sur un autre « \&c. » empâté. La caractéristique $\delta$ est un « d » à boucle, distinct du $d$ droit. « on a lors » : ainsi, pour « alors ». Le crochet est sur la feuille, comme aux vues précédentes.}

Pag 168 et suivantes Memoire de La Grange sur la methode des
\emph{variations}.
\note{« variations » est souligné. Un trait horizontal sépare ce titre de ce qui précède.}

On suppose \quad \uncertain{$\psi$} $= P\delta\varphi + Q\delta x + R\delta y + S\delta z + \&c.$
\[ \Pi = P'\delta\varphi + P''d\delta\varphi + P'''d^{2}\delta\varphi + \&c. \]
\[ + Q'\delta x + Q''d\delta x + Q'''d^{2}\delta x + \&c. \]
\[ + R'\delta y + R''d\delta y + R'''d^{2}\delta y + \&c. \]
\[ + \&c. \]
\[ \int\psi + \Pi = a \text{ une constante} \]
\note{La lettre que l'on donne pour $\psi$ est un « v » barré, marqué d'un point dans les deux lignes suivantes. $\Pi$ est tracé comme deux jambages sous une barre. Au-dessus du « O » de « On », une petite marque d'encre en forme de « E », hors du texte.}

On prouve par le calcul que si on change $\delta$ en $d$ dans l'expression
de $\Pi$ on aura toujours $\Pi = 0$ et par conséquent $\int\psi = \text{constante}$
\[ \psi = 0 \]
\note{« l'expression » est récrit sur un premier mot, peut-être « la ». $\psi = 0$ est écrit sous « $\int\psi = $ constante ».}

[ainsi dans Condorcet ce sont les quantités sous le signe qui deviennent nulle tandis que le
le même changement de $\delta$ en $d$ rend nulle les quantités hors le signe dans le present
mémoire; sans doute cette difference vient de changement dans les hypotheses
mais pour en bien juger il faudroit avoir un mémoire précédent sous les yeux]
\note{Ce paragraphe entre crochets est d'une écriture plus petite. « le le » à la coupure de ligne, « nulle » au singulier : ainsi sur la feuille. Le bas de la page est blanc.}

\page{747}

\note{Verso du feuillet 371 (vue 746) ; aucun numéro écrit de ce côté. La page poursuit les extraits des volumes de Turin : Lagrange encore, sur le mouvement d'un corps attiré vers des centres fixes.}

[On trouve Pag. 188 jusqu'a \uncertain{243} deux memoires de la Grange
qu'il a analysés en y fesant quelque transformation et addition
dans la section de la mechanique ou il traite
Du mouvement d'un corps attiré vers un ou plusieurs centres]
\note{« la Grange » : le « n » est ajouté au-dessus de la ligne. Le deuxième chiffre de « 243 » a la forme bouclée du 4 de la vue 746. « la mechanique » renvoie, semble-t-il, à la Mécanique analytique de Lagrange, citée plus bas par sa page ; c'est une inférence.}

[Dans le cas ou le corps est attiré par deux centres fixes]
\marginal{Pag. 210.}
Si l'un des centres des forces s'eloignoit a l'infini, alors la force
tendente a ce centre infiniment distant deviendroit uniforme, et
agiroit suivant des lignes parallèles; ainsi on auroit le cas d'un
corps attiré vers un centre fixe par une force reciproquement
proportionnelle aux quarrés des distances, et poussé en même
tems par une force devaleur et de direction constante.
\note{« Pag. 210. » est écrit dans la marge de gauche, en face de « Si l'un des centres ». Un crochet ouvrant, isolé, se voit aussi dans la marge en face de « proportionnelle ».}

[Pour appliquer les formules de la méchanique (\emph{Pag. 271}) au cas present
il est visible que l'on doit faire $b = \infty$ et $r$ ou $q$ egalement $\infty$]

2\textsuperscript{ème} Memoire \emph{Pag. 218}
\note{Ce titre est souligné d'un long trait.}

Proposons nous les équations $\frac{d^{2}x}{2dt^{2}} + M = 0$, \quad $\frac{d^{2}y}{2dt^{2}} + N = 0$, $M$ et $N$ étants
des fonctions de $x$, $y$; et voyons quelles sont les conditions de l'intégrabilité
de ces équations.
Nous supposerons que ces deux équations étants multipliées
l'une et l'autre par des quantités convenables, et ensuite ajoutées
ensemble forment une équation intégrable. Cette supposition est
peut-être la seule qui convienne a la nature des équations proposées,
et elle est d'ailleurs la plus naturelle, et en même tems la plus
générale qu'on puisse faire. Or comme les deux membres $M$ et
$N$ sont des quantités finies, il est visible que les multiplicateurs
ne sauroient être que des quantités differentielles du premier ordre.
Prenons donc $m\,dx + n\,dy$, et $\mu\,dx + \nu\,dy$ pour les multipicateurs dont il
s'agit, $m$, $n$, $\mu$ et $\nu$ étant des fonctions de $x$ et $y$ et nous aurons l'équation
\[ \frac{m\,dx\,d^{2}x + n\,dy\,d^{2}x + \mu\,dx\,d^{2}y + \nu\,dy\,d^{2}y}{2dt^{2}} \pm (mM + \mu N)\,dx + (nM + \nu N)\,dy = 0 \]
laquelles ne sauroit être intégrable amoins que les deux parties (savoir
celle qui contient les differences secondes, et celle qui ne contient que
les differences premières) ne soient intégrable chacune en particulier.
\note{Devant « Nous supposerons », une grande croix d'encre à boucle, en marge ; peut-être une majuscule manquée et reprise. « étants », « multipicateurs », « laquelles », « intégrable » : ainsi. Les deux derniers multiplicateurs sont tracés comme « u » (ou « c ») et « v » ; on les donne $\mu$ et $\nu$, choix du transcripteur. Le signe devant $(mM + \mu N)$ est un $+$ posé sur un trait, lu $\pm$.}

Or il n'est pas difficile de voir que l'intégrale de la première ne
peut être que de cette forme $\frac{A\,dx^{2} + B\,dx\,dy + C\,dy^{2}}{2dt^{2}}$, $A$, $B$ $C$,
étant des fonctions de $x$ et $y$; donc comparant la differentielle
de la quantité $A\,dx^{2} + B\,dx\,dy + C\,dy^{2}$ qui est
\[ 2A\,dx\,d^{2}x + B(dx\,d^{2}y + dy\,d^{2}x) + 2C\,dy\,d^{2}y + \left(\frac{dA}{dy} + \frac{dB}{dx}\right)dx\,dy \]
\[ + \left(\frac{dB}{dy} + \frac{dC}{dx}\right)dy^{2}\,dx + \frac{d.C}{dy}\,dy^{3}, \quad \text{avec la quantité} \]
\[ m\,dx\,d^{2}x + \mu\,dx\,d^{2}y + n\,dy\,d^{2}x + \nu\,dy\,d^{2}y, \quad \text{on aura} \]
\note{La différentielle recopiée est incomplète : il y manque les termes $\frac{dA}{dx}dx^{3}$ et, au lieu de $\left(\frac{dA}{dy} + \frac{dB}{dx}\right)dx\,dy$, on attendrait $dx^{2}\,dy$ ; transcrite telle quelle. La phrase s'arrête sur « on aura » au bas du feuillet, et la suite est à la vue 748.}

\page{748}

\note{En haut à droite, à l'encre, « 372 », au bout de la première ligne. La page commence par le résultat de la comparaison annoncée au bas de la vue 747 (« on aura »). Le cercle pâle qu'on devine au milieu de la page est le cachet de la vue 749 vu par transparence.}

\[ 2A = m,\ B = n = \mu,\ 2C = \nu,\ \frac{dA}{dx} = 0,\ \frac{dC}{dy} = 0,\ \frac{dA}{dy} + \frac{dB}{dx} = 0,\ \frac{dB}{dy} + \frac{dC}{dx} = 0. \]
Les équations $\frac{dA}{dx} = 0$ et $\frac{dC}{dy} = 0$ donnent d'abord $A = \text{fonc.}\,y$, $C = \text{fonc.}\,x$.
[en effet $\frac{dA}{dx}$ étant le coëfficient \add{\uncertain{de dans}} de la differentielle de $A$ prise par rapport
a $x$, il en resulte evidemment que $A$ \add{ne} contient \add{pas} $x$, \struck{puisque sans cette condition}
\struck{$A$ ne pourroit être differentiée par rapport a cette variable}
\add{car si $A$ etoit fonction de $x$ la differentielle partielle prise parrapport a $x$ ne seroit pas nulle. Mais $A$ \struck{ab} a été d'abord supposée fonction de $y$ et $x$ donc puis qu'elle ne contient pas $x$ elle est simplement fonction de $y$. le même raisonnement p. $C$.}
Dans le calcul des fonctions on auroit eu $A' = 0$ et la manière deplacer
l'accent auroit indiqué par rapport a quelle variable la fonction prime auroit été
prise, cequi revient absolument au même]
\note{Dans la première ligne, « $n =$ » est récrit sur un autre signe. Les lettres $y$ et $x$ de « fonc.\ $y$ », « fonc.\ $x$ » sont chacune récrites sur l'autre. Au-dessus de « de la differentielle », en plus petit, « de dans » (ou « dedans »), sans marque d'insertion. Le passage de « car si A etoit » à « p. C. » est écrit en très petit au-dessus et au-dessous des lignes barrées, qu'il remplace ; « p. C. » vaut sans doute « pour C ». Deux petits traits dans la marge de gauche, en face de « en effet ».}

Ensuite les équations $\frac{dA}{dy} + \frac{dB}{dx} = 0$, et $\frac{dB}{dy} + \frac{dC}{dx} = 0$ donnent $\frac{d^{2}A}{dy^{2}} = \frac{d^{2}C}{dx^{2}}$
d'ou il aisé de conclure que les quantités $A$ et $C$ ne peuvent que dela forme
suivante $A = a + by + cy^{2}$, $C = g + hx + cx^{2}$, $a$, $b$, $c$, $g$, $h$ étant des constantes.
[on voit effectivement que les valeurs de $A$ et $C$ ne peuvent pas contenir des puissances
des variables, dont elles sont respectivement fonction, plus elevées que le second
degré il est de plus visible que cette puissance seconde \struck{de} ne peut être
multipliée que par la même constante dans l'une et l'autre fontion, car
cette dernière condition est également necessaire pour que l'équation
$\frac{d^{2}A}{dy^{2}} = \frac{d^{2}C}{dx^{2}}$ ait lieu. Soit donc $\frac{d^{2}A}{dy^{2}} = 2c$ ou $\frac{d^{2}A}{dy} = 2c\,dy$ on a en intégrant
\quad $\frac{dA}{dy} = 2cy + b$, $dA = 2cy\,dy + b\,dy$
\quad $A = cy^{2} + by + a$
en fesant la même opération pour $C$
on trouve la valeur indiquée \struck{\ill{}}. on voit que
les constantes $b$, $a$, peuvent être differentes de $g$ et $h$ parceque ces
lettres disparoissent \struck{\ill{}} en differentiant deux fois]
\note{« d'ou il aisé » : ainsi, sans « est » ; « ne peuvent que dela forme » : ainsi, sans « être ». Les trois lignes du calcul de $A$ sont écrites à droite de la page, en face de « en fesant \ldots{} on voit que ». Après « indiquée », un mot court noyé sous une tache ; après « disparoissent », un mot de deux lettres barré.}

Maintenant l'équation $\frac{dA}{dy} + \frac{dB}{dx} = 0$, $\frac{dB}{dx} = -b - 2cy$ d'ou en intégrant $B = -(b + 2cy)x + \text{fonc.}\,y$
[on ajoute fonc.\ $y$ parceque $B$ ayant été supposée fonction de $x$ et $y$ on doit ajouter une fonc.\ de $y$ a l'intégrale
de la differentielle prise pour $x$] mais l'autre équation $\frac{dB}{dy} + \frac{dC}{dx} = 0$ donnera demême
$B = -(h + 2cx)y + \text{fonc.}\,x$ ensorte qu'on aura $B = k - bx - hy - 2cxy$, $k$ étant une constante]
\note{Le $y$ de « fonc.\ $y$ », en fin de première ligne, est empâté. Le crochet fermant de la dernière ligne n'a pas de crochet ouvrant correspondant après le premier, déjà fermé.}

[car engénéral $\frac{dB}{dx}dx + \frac{dB}{dy}dy = dB$ mais ici on a $\frac{dB}{dx}dx + \frac{dB}{dy}dy = -(b + 2cy)\,dx - (h + 2cx)\,dy$
et l'on tire de la première valeur de $B$, $dB = -(b + 2cy)\,dx - 2cx\,dy + \text{fonc}\,y'\,dy$ puis de
la seconde \ldots\ldots \quad $dB = -(h + 2cx)\,dy - 2cy\,dx + \text{fonc.}\,x'\,dx$
et comparant ces valeurs de $dB$ avec celle de la differentielle partielle
on a $-h\,dy = \text{fonc}\,y'\,dy$, \quad $-b\,dx = \text{fonc.}\,x'\,dx$
\quad $-hy = \text{fonc.}\,y + \text{constante}$, \quad $-bx = \text{fonc.}\,x + \text{constante}$ \struck{si on}
si donc on met pour fonc.\ $x$ et fonc.\ $y$ leurs valeurs dans celle $B$ on aura
l'équation indiqué]
\note{Dans « $-h\,dy = \text{fonc}\,y'\,dy$ », le $y$ et le $dy$ sont récrits ; dans « $-bx$ », le $x$ est récrit sur un autre signe. « si on », en bout de ligne, est barré et repris à la ligne suivante.}

Ayant déterminé ainsi les quantités $A$, $B$, \struck{et} $C$, on aura pour les multiplicateurs des
équations proposées, les deux quantités $2A\,dx + B\,dy$, et $B\,dx + 2C\,dy$; et il ne
restera plus qu'a rendre intégrables la formule $(2AM + BN)\,dx + (BM + 2CN)\,dy$,
car alors, en nommant $Z$ l'intégrale de cette formule, on aura l'équation
dupremier ordre. \quad $\frac{A\,dx^{2} + B\,dx\,dy + C\,dy^{2}}{2dt^{2}} + Z = \text{une constante}$ \struck{$= a$}
Or pour qu'une quantité telle que $(2AM + BN)\,dx + (BM + 2CN)\,dy$ soit une
differentielle exacte il faut que l'on ait $\frac{d.(2AM + BN)}{dy} = \frac{d.(BM + 2CN)}{dx}$, c'est
l'équation de condition qui doit avoir lieu pour que les équations proposées admettent une intégrale du premier
ordre.
\note{« une constante » : « une » est écrit au-dessus de « constante », qui remplace un « $= a$ » barré après $Z$, et un second « $= a$ » barré en bout de ligne ; les deux ratures sont réunies ici. La dernière ligne, serrée, finit au bord droit du feuillet (« ordre » est écrit au-dessus de la fin de la ligne). Le texte s'arrête au bas de la page et reprend en tête de la vue 749.}

\page{749}

\note{Verso du feuillet 372 (vue 748) ; aucun numéro écrit de ce côté. Dans la moitié haute, à gauche, le cachet rond « BIBLIOTHEQUE ROYALE » à la couronne, posé sur les équations. La page continue d'abord le second mémoire de Lagrange de la vue 748, puis passe, sous un trait épais, à un extrait sur l'équation $x^{2} - ay^{2} = 1$.}

Puisque $(2AM + BN)\,dx + (BM + 2CN)\,dy = dZ$, on aura $2AM + BN = \frac{dZ}{dx}$, et
$BM + 2CN = \frac{dZ}{dy}$ et parconséquent $M = \frac{2C\frac{dZ}{dx} - B\frac{dZ}{dy}}{4AC - B^{2}}$, $N = \frac{2A\frac{dZ}{dy} - B\frac{dZ}{dx}}{4AC - B^{2}}$.
où l'on pourra prendre pour $Z$ une fonction quelconque de $x$, et $y$.
Si on fait encore $4AC - B^{2} = D$, et $Z = D^{2}V$, ($V$ etant une fonction
quelconque de $x$ et $y$) on aura \quad $M = 2V\left(2C\frac{dD}{dx} - B\frac{dD}{dy}\right) + D\left(2C\frac{dV}{dx} - B\frac{dV}{dy}\right)$
\[ N = 2V\left(2A\frac{dD}{dy} - B\frac{dD}{dx}\right) + D\left(2A\frac{dV}{dy} - B\frac{dV}{dx}\right) \]
or a cause de $\frac{dA}{dx} = 0$, $\frac{dC}{dy} = 0$, $\frac{dB}{dx} = -\frac{dA}{dy}$ et $\frac{dB}{dy} = -\frac{dC}{dx}$, on trouvera
\[ \frac{dD}{dx} = 2\left(2A\frac{dC}{dx} + B\frac{dA}{dy}\right)\ \frac{dD}{dy} = 2\left(2C\frac{dA}{dy} + B\frac{dC}{dx}\right);\ \text{donc on aura} \]
\[ M = D\left(4V\frac{dC}{dx} + 2C\frac{dV}{dx} - B\frac{dV}{dy}\right), \]
\[ N = D\left(4V\frac{dA}{dy} + 2A\frac{dV}{dy} - B\frac{dV}{dx}\right), \]
Ainsi toutes les fois que les quantités $M$ et $N$ pourront se ramener a
cette forme les équations $\frac{d^{2}x}{dt^{2}} + M = 0$, et $\frac{d^{2}y}{dt^{2}} + N = 0$ auront pour
intégrale \quad $\frac{A\,dx^{2} + B\,dx\,dy + C\,dy^{2}}{2dt^{2}} + D^{2}V = a$ une constante
\note{« $\frac{dC}{dy} = 0$ » et le début de « $\frac{dB}{dx}$ » sont sous le cachet et se lisent à travers. Dans la seconde expression de $\frac{dD}{dy}$, le « 2 » est récrit sur un autre chiffre. Ici les équations sont écrites sans le « 2 » au dénominateur ($\frac{d^{2}x}{dt^{2}}$), à la différence de la vue 747. Dans la dernière formule, le $dy$ de $C\,dy^{2}$ est récrit.}

[Dans le mémoire cité (\emph{P. 37}) dela Théorie des nombres on trouve (P. \uncertain{46} T. \uncertain{4} de Turin)
\note{Le premier chiffre de « 46 » et le « 4 » de « T. 4 » ont la forme bouclée du 4 des vues 743 et 746. « la Théorie des nombres » renvoie à un ouvrage qui cite, à sa page 37, le mémoire de Lagrange du tome 4 de Turin ; le crochet ouvert n'est pas refermé sur cette ligne.}

Lemme. Le produit de ces deux quantités $x^{2} - ay^{2}$ et $x'^{2} - ay'^{2}$ est $(xx' \pm ayy')^{2} - a(xy' \pm yx')^{2}$;
car \quad $(x^{2} - ay^{2})(x'^{2} - ay'^{2}) = x^{2}x'^{2} + a^{2}y^{2}y'^{2} - ay^{2}x'^{2} - ax^{2}y'^{2}$.
\[ = x^{2}x'^{2} \pm 2axx'yy' + a^{2}y^{2}y'^{2} - ax^{2}y'^{2} \mp 2axyx'y' - ay^{2}x'^{2} \]
\[ = (xx' \pm ayy')^{2} - a(xy' \pm yx')^{2} \]
[et P. 60 du même mémoire]

Or comme le quarré le cube \struck{\&c.} et en général toute puissance d'une quantité de cette forme
$x^{2} - ay^{2}$ est toujours aussi de la même forme il s'en suit qu'en élevant l'équation
$1 = x^{2} - ay^{2}$ à une puissance quelconque, on aura une infinité d'autres équations semblables
de sorte qu'ayant trouvé par quelque méthode que ce soit, une seule solution du
problème, on pourra par son moyen en trouver d'autres a l'infini.
Pour renfermer toutes ces solutions dans une formule générale, supposons que
$p$ et $q$ soient les valeurs trouvées de $x$ et $y$ en sorte que l'on ait $1 = p^{2} - aq^{2}$; en élevant les
deux membres de cette équation a une puissance quelconque, $m$, on aura $1 = (p^{2} - aq^{2})^{m}$,
équation qu'il s'agit de réduire a la forme de celle-ci $1 = x^{2} - ay^{2}$.
Pour cela je remarque que $p^{2} - aq^{2} = (p + q\sqrt{a})(p - q\sqrt{a})$, de sorte que l'on aura
$(p^{2} - aq^{2})^{m} = (p + q\sqrt{a})^{m}(p - q\sqrt{a})^{m}$. or, $(p + q\sqrt{a})^{m} = p^{m} + mp^{m-1}q\sqrt{a} + \frac{m(m-1)}{2}.p^{m-2}q^{2}a + \&c.$
donc si on fait $x = p^{m} + \frac{m(m-1)}{2}p^{m-2}q^{2}a + \&c.$ \quad $y = mp^{m-1}q + \frac{m(m-1)(m-2)}{2.3}.p^{m-3}q^{3}a + \&c.$
on aura $(p + q\sqrt{a})^{m} = x + y\sqrt{a}$; et prenant le radical $\sqrt{a}$ en $-$, on aura demême
$(p - q\sqrt{a})^{m} = x - y\sqrt{a}$ donc $(p^{2} - aq^{2})^{m} = (x + y\sqrt{a})(x - y\sqrt{a}) = x^{2} - ay^{2}$; de sorte que l'on aura
engénéral $x^{2} - ay^{2} = 1$, en prenant pour $m$ un nombre entier et positif.
Au reste les équations $(p + q\sqrt{a})^{m} = x + y\sqrt{a}$, et $(p - q\sqrt{a})^{m} = x - y\sqrt{a}$ donneront.
\[ x = \frac{(p + q\sqrt{a})^{m} + (p - q\sqrt{a})^{m}}{2}, \quad y = \frac{(p + q\sqrt{a})^{m} - (p - q\sqrt{a})^{m}}{2\sqrt{a}} \]
expressions qui reviennent au même que les précédentes; mais qui ont l'avantage d'être sous une forme
finie, ainsi prenant successivement pour $m$ \struck{a tous} les nombres \struck{\ill{}} \add{naturels} on aura une infinité de solutions.
\note{Après « le cube », un « \&c. » barré. « $m$ » est écrit au-dessus de « pour », sur « a tous » barré ; « naturels », au-dessus de la ligne, remplace un mot barré qu'on ne lit pas ; « on aura » est récrit sur la rature.}

[La Grange demontre ensuite \emph{P. 65}] que si $p$ et \struck{\ill{}} $q$ sont les plus petites de $x$ et $y$ qui satisfassent
a l'équation $x^{2} - ay^{2} = 1$, toutes les autres valeurs de $x$ et de $y$ sont necessairement renfermées
dans les formules générales. Pour cela il prouve que si $p$ et $q$ sont les plus petites valeurs et
$p'$ et $q'$ \struck{soient} les plus petites après celle on aura \quad $p' = \frac{(p + q\sqrt{a})^{2} + (p - q\sqrt{a})^{2}}{2}$, $q' = \frac{(p + q\sqrt{a})^{2} + (p - q\sqrt{a})^{2}}{2\sqrt{a}}$.
[il est visible que si on fait $m = 1$ les deux membres des équations deviennent identiques cequi \struck{s'accorde} doit
être en effet pour que la proposition soit vraie]
\note{Devant « $q$ », un mot bref barré. Dans $q'$, le numérateur porte « $+$ » entre les deux carrés, là où la formule générale de $y$ donne « $-$ » ; laissé tel. Ce paragraphe, entre crochets, est d'une écriture plus petite et serrée. Le bas de la page est blanc ; c'est la dernière page écrite du volume.}

\end{document}
